\section{引言：阅读 K2.5 的三条主线}
\begin{figure}[H]
\centering
\includegraphics[width=0.78\textwidth]{cover.jpg}
\caption{K2.5 在发布片花中的 ``Visual Agency Intelligence'' 意象，强调视觉代理能力的起点。 \protect\footnotemark}
\end{figure}
\footnotetext{视频画面时间区间：00:00:05--00:00:13，讲师开场时点出的 K2.5 技术路线。}
月之暗面在 2026 年 1 月 25/26 日正式亮相 K2.5，讲师在开场就说：\textit{``Dive into Kimi K2.5 这篇庞杂的文章''}，并把这一整期课定位为从数据、算法到基础设施的整体串讲。两天后发布的 Gemini 3 Flash agent vision 版本也在这个时间线上被点名，目的是说明原生多模态与视觉代理不是孤立事件。

\begin{knowledgebox}{三条阅读线索}
\begin{itemize}
    \item \textbf{Data}：追问每个新数据域是如何定义、采样、合成，新的 domain 就意味着新的能力。
    \item \textbf{Algorithm}：谁负责评判模型回答？GRM 与多阶段训练如何协同稳住 open-ended 任务。
    \item \textbf{Infra}：如何用 Critical Steps、并行 agent swarm 和评估流水线共同管理 compute 与 latency。
\end{itemize}
\end{knowledgebox}

\subsection{本章小结}
K2.5 的导读就是从 ``Read the tech report'' 出发，明确我们要用数据、算法、架构三个视角追踪 Visual Agency Intelligence 的实现路径，并把这份讲稿当成技术报告的注释版。
